\documentclass[12pt, a4paper, toc=sectionentrywithdots]{scrartcl}

% ------------------------------------------------------------------------------
% 1. KONFIGURATION & PAKETE
% ------------------------------------------------------------------------------
\usepackage{../Organisatorisch/JKU_SVP} % für Design und Befehle
% --- ALLGEMEINE EINSTELLUNGEN FÜR ALLE PROTOKOLLE ---

% Tragt hier das aktuelle Semester ein
\newcommand{\Semester}{Sommersemester 2026}

% Tragt hier eure Gruppennummer ein
\newcommand{\GruppenNummer}{0} 

% Tragt hier eure Namen ein (getrennt durch Komma)
\newcommand{\GruppenMitglieder}{Vorname Nachname, Vorname Nachname}

% Definiere die Farbe aller Überschriften (RGB-Format)
\definecolor{SectionColor}{RGB}{0,40,210}

\colorlet{SubSubSectionColor}{SectionColor!65!white} % Farbe der Unterunter-Kapitel

% Sollen die Überschriften mit oder ohne Serifen sein?
    % \SerifHeadingstrue... Überschriften in Serifenschrift.
    % \SerifHeadingsfalse... Überschrift in serifenloser Schrift.
\SerifHeadingstrue



% =========== GLOBAL: Hier eigene Pakete ergänzen (\usepackage{...}) ===========
%
% Die hier eingefügten Pakete und Befehle sind in ALLEN Protokollen aktiv und
% müssen somit nur einmal eingegeben werden.
% ------------------------------------------------------------------------------
            % Eistellungen.tex initial einrichten

\addbibresource{Quellen.bib}

% ==============================================================================
%
%                     /$$$$$$      /$$    /$$     /$$$$$$$ 
%                    /$$__  $$    | $$   | $$    | $$__  $$
%                   | $$  \__/    | $$   | $$    | $$  \ $$
%                   |  $$$$$$     |  $$ / $$/    | $$$$$$$/
%                    \____  $$     \  $$ $$/     | $$____/ 
%                    /$$  \ $$      \  $$$/      | $$      
%                   |  $$$$$$/       \  $/       | $$      
%                    \______/         \_/        |__/      
%                                                                  
%                SCHULVERSUCHSPRAKTIKUM PHYSIK | LATEX VORLAGE
%
% ==============================================================================


% ------------------------------------------------------------------------------
% ================ Hier die Infos zu diesem Protokoll eintragen ================
% ------------------------------------------------------------------------------
\praktikumstag  {17. März 2026} % Tag, an dem der Versuch durchgeführt wurde
\versuchstitel  {XX - Thema} % "<Kurzbez.> - <Name des Themas>"


% ========== Hier optional eigene Pakete ergänzen (\usepackage{...}) ===========
%
% Pakete und Befehle die in allen Protokollen verwendet werden, sollten in
% Einstellungen.tex gegeben werden.
% ------------------------------------------------------------------------------



% ==============================================================================
% VERLINKUNGEN & PDF-METADATEN (Hyperref)
% ------------------------------------------------------------------------------
\usepackage[
    colorlinks=true,  
    linkcolor=black,   
    citecolor=black,   
    urlcolor=black,    
    bookmarksnumbered=true
]{hyperref}
 
\usepackage[ngerman]{cleveref}
% ------------------------------------------------------------------------------

% ------------------------------------------------------------------------------
% ==============================================================================
% %%%%%%%%%%%%%%%%%%%%%%%%%%%%  DOKUMENT BEGINN  %%%%%%%%%%%%%%%%%%%%%%%%%%%%%%%
% ==============================================================================
% ------------------------------------------------------------------------------
\begin{document}

% ------------------------------- TITELSEITE  ----------------------------------
\maketitlepageSVP

% =========================== HIER STARTET DER INHALT ==========================
%
% Der Inhalt wird grundsätzlich in die einzelnen Kapitel .tex gegeben, kann
% aber auch hier direkt statt \include{} geschrieben werden.
%    Für Hauptüberschriften wird \addsec{Versuch 1: ...} verwendet.
%    Für Unterüberschirften wird \subsection*{...} verwendet.
% ------------------------------------------------------------------------------

\addsec{Didaktische Basisinformationen zu allen Versuchen}


\subsection*{Alltagsbezug – Wozu brauche ich denn das?}
% -------------------------------------------
%                ALLTAGSBEZUG
% -------------------------------------------




\subsection*{Engage-Idee zum Thema}
% -------------------------------------------
%               ENGAGE-IDEE
% -------------------------------------------



\subsection*{Schülervorstellungen}
% -------------------------------------------
%            SCHÜLERVORSTELLUNGEN
% -------------------------------------------




\include{./Kapitel/Versuch1}


\include{./Kapitel/Versuch2}


%%%%%%%%%%%%%%%%%%%%%%%%%%%%%%%%%%%%%%%%%%%%%
% >>>>>>>>>> TITEL DES VERSUCHS <<<<<<<<<<<<<
%%%%%%%%%%%%%%%%%%%%%%%%%%%%%%%%%%%%%%%%%%%%%
\addsec{Versuch 3: [Aussagekräftiger Name des Versuchs]}


\subsection*{Key-Ideas}
% -------------------------------------------
%                 KEY-IDEAS
% -------------------------------------------




\subsection*{Ziel des Versuchs}
% -------------------------------------------
%                   ZIELE
% -------------------------------------------




\subsection*{Begründung für E}
% -------------------------------------------
%             BEGRÜNDUNG FÜR E
% -------------------------------------------




\subsection*{Durchführung}
% -------------------------------------------
%               DURCHFÜHRUNG
% -------------------------------------------

\subsubsection*{Material}
% -------------------------------------------
%                  MATERIAL
% -------------------------------------------




\subsubsection*{Aufbau und Skizze}
% -------------------------------------------
%              AUFBAU und SKIZZE
% -------------------------------------------




\subsubsection*{Ablauf}
% -------------------------------------------
%                   ABLAUF
% -------------------------------------------




\subsubsection*{Ergebnisse und kurze Erklärung}
% -------------------------------------------
%                ERGEBNISSE
% -------------------------------------------




\subsubsection*{Tipps \& Tricks}
% -------------------------------------------
%               TIPPS & TRICKS
% -------------------------------------------


% -------  Literatur  --------
\printbibliography[title={Literatur}, heading=bibintoc]


\end{document}
